\documentclass{article}
\usepackage{lmodern}
\usepackage{amsmath}
\usepackage{listings}
\usepackage{fullpage}
\usepackage{parskip}
\usepackage{xcolor}
\lstset{
	basicstyle=\ttfamily,
	columns=fullflexible,
	frame=single,
	breaklines=true,
	postbreak=\mbox{\textcolor{red}{$\hookrightarrow$}\space},
}
\begin{document}
\title{Experiment `p\_6\_2\_f\_move\_even\_to\_front\_array' Results}
\author{\today}
\date{}
\maketitle
\textbf{Experiment outcome: }SUCCESS
\\\textbf{Bad responses: }0
\\\textbf{Responses containing}~\texttt{assume}~\textbf{: }0
\\\textbf{Responses containing}~\texttt{expect}~\textbf{: }0
\\\textbf{Resolution attempts: }5
\\\textbf{Hard fails (resolution): }3
\\\textbf{Soft fails (resolution): }0
\\\textbf{Verification attempts: }2
\section*{Problem Specification}
\textbf{Problem name: }p\_6\_2\_f\_move\_even\_to\_front\_array
\\\textbf{Natural language statement: }Write a method to move all even elements in an array to the front, otherwise preserving the order of the elements.
\\\textbf{Method signature: }p\_6\_2\_f\_move\_even\_to\_front\_array(inputs: array<int>) returns (result: array<int>)
\subsection*{Ensures}
\begin{itemize}
\item \texttt{result[..] == evens(inputs[..]) + odds(inputs[..])}
\end{itemize}
\subsection*{Functional Code Given}
\begin{lstlisting}
function evens(s: seq<int>): seq<int>
{
  if |s| == 0 then []
  else if s[0] % 2 == 0
       then [s[0]] + evens(s[1..])
       else evens(s[1..])
}

function odds(s: seq<int>): seq<int>
{
  if |s| == 0 then []
  else if s[0] % 2 != 0
       then [s[0]] + odds(s[1..])
       else odds(s[1..])
}
\end{lstlisting}
\clearpage
\section*{GenAI interactions}
Below you will find all interactions between the `user' (program) and the `assistant' (OpenAI).
\subsection*{Program $\rightarrow$ GenAI}
\begin{lstlisting}
You are given the following task to perform in Dafny:

Write a method to move all even elements in an array to the front, otherwise preserving the order of the elements.

The signature should be:

p_6_2_f_move_even_to_front_array(inputs: array<int>) returns (result: array<int>)

The method should respect the following contract:

ensures result[..] == evens(inputs[..]) + odds(inputs[..])

The contract uses the following dafny code:

function evens(s: seq<int>): seq<int>
{
  if |s| == 0 then []
  else if s[0] % 2 == 0
       then [s[0]] + evens(s[1..])
       else evens(s[1..])
}

function odds(s: seq<int>): seq<int>
{
  if |s| == 0 then []
  else if s[0] % 2 != 0
       then [s[0]] + odds(s[1..])
       else odds(s[1..])
}

These function(s) must not be used in your implementation of the method.

Produce and show only the Dafny body of this method, including the curly braces that surround it. Do not show the signature nor contract. You must not use 'assume' nor 'expect' anywhere in your code.
\end{lstlisting}
\subsection*{GenAI $\rightarrow$ Program}
\textbf{System fingerprint: }null
\\\textbf{ID: }chatcmpl-Dn0z96CwfVTuMUtXa3V59GaUrbslv
\\\textbf{Model: }gpt-5.5-2026-04-23
\\\textbf{Created at: }1780574231
\\\textbf{Finish reason: }stop
\\\textbf{Usage: }588 tokens in, and 8730 tokens out
\begin{lstlisting}
{
  var n := inputs.Length;
  var es: seq<int> := [];
  var os: seq<int> := [];
  var i := n;

  assert inputs[n..] == [];
  assert evens(inputs[n..]) == [];
  assert odds(inputs[n..]) == [];

  while 0 < i
    invariant 0 <= i <= n
    invariant es == evens(inputs[i..])
    invariant os == odds(inputs[i..])
  {
    var oldI := i;
    var k := i - 1;
    var x := inputs[k];
    ghost var tail := inputs[oldI..];
    ghost var whole := inputs[k..];

    assert 0 <= k < n;
    assert oldI == k + 1;
    assert inputs[k..oldI] == [x];
    assert whole == inputs[k..oldI] + tail;
    assert whole == [x] + tail;
    assert |whole| != 0;
    assert whole[0] == x;
    assert whole[1..] == tail;

    if x % 2 == 0 {
      es := [x] + es;
      assert evens(whole) == [x] + evens(tail);
      assert odds(whole) == odds(tail);
      assert es == evens(inputs[k..]);
      assert os == odds(inputs[k..]);
    } else {
      os := [x] + os;
      assert evens(whole) == evens(tail);
      assert odds(whole) == [x] + odds(tail);
      assert es == evens(inputs[k..]);
      assert os == odds(inputs[k..]);
    }

    i := k;
  }

  assert i == 0;
  assert inputs[i..] == inputs[..];
  assert es == evens(inputs[..]);
  assert os == odds(inputs[..]);

  var combined := es + os;
  result := new int[|combined|](j => combined[j]);

  assert result.Length == |combined|;
  assert forall j :: 0 <= j < |combined| ==> result[j] == combined[j];
  assert |result[..]| == |combined|;
  assert forall j | 0 <= j < |combined| ensures result[..][j] == combined[j] {
    assert result[..][j] == result[j];
  }
  assert result[..] == combined;
  assert combined == evens(inputs[..]) + odds(inputs[..]);
}
\end{lstlisting}
\subsection*{Program $\rightarrow$ GenAI}
\begin{lstlisting}
When using dafny resolve, the below error is emitted and resolve fails:

problem_data/p_6_2_f_move_even_to_front_array/program.dfy(82,0): Error: invalid UnaryExpression
   |
82 | }
   | ^

problem_data/p_6_2_f_move_even_to_front_array/program.dfy(81,57): Error: expected either 'by' or a semicolon following the assert expression
   |
81 |   assert combined == evens(inputs[..]) + odds(inputs[..]);
   |                                                          ^

2 parse errors detected in program.dfy

Correct the error in code by altering only the method body. Produce and show only the Dafny body, including the curly braces that surround it. Do not show the signature nor contract. You must not use 'assume' nor 'expect' anywhere in your code.
\end{lstlisting}
\subsection*{GenAI $\rightarrow$ Program}
\textbf{System fingerprint: }null
\\\textbf{ID: }chatcmpl-Dn129NrLlM9VEnjUwh9KdqdPOWnGe
\\\textbf{Model: }gpt-5.5-2026-04-23
\\\textbf{Created at: }1780574417
\\\textbf{Finish reason: }stop
\\\textbf{Usage: }1275 tokens in, and 1036 tokens out
\begin{lstlisting}
{
  var n := inputs.Length;
  var es: seq<int> := [];
  var os: seq<int> := [];
  var i := n;

  assert inputs[n..] == [];
  assert evens(inputs[n..]) == [];
  assert odds(inputs[n..]) == [];

  while 0 < i
    invariant 0 <= i <= n
    invariant es == evens(inputs[i..])
    invariant os == odds(inputs[i..])
  {
    var oldI := i;
    var k := i - 1;
    var x := inputs[k];
    ghost var tail := inputs[oldI..];
    ghost var whole := inputs[k..];

    assert 0 <= k < n;
    assert oldI == k + 1;
    assert inputs[k..oldI] == [x];
    assert whole == inputs[k..oldI] + tail;
    assert whole == [x] + tail;
    assert |whole| != 0;
    assert whole[0] == x;
    assert whole[1..] == tail;

    if x % 2 == 0 {
      es := [x] + es;
      assert evens(whole) == [x] + evens(tail);
      assert odds(whole) == odds(tail);
      assert es == evens(inputs[k..]);
      assert os == odds(inputs[k..]);
    } else {
      os := [x] + os;
      assert evens(whole) == evens(tail);
      assert odds(whole) == [x] + odds(tail);
      assert es == evens(inputs[k..]);
      assert os == odds(inputs[k..]);
    }

    i := k;
  }

  assert i == 0;
  assert inputs[i..] == inputs[..];
  assert es == evens(inputs[..]);
  assert os == odds(inputs[..]);

  var combined := es + os;
  result := new int[|combined|](j => combined[j]);

  assert result.Length == |combined|;
  assert forall j :: 0 <= j < |combined| ==> result[j] == combined[j];
  assert |result[..]| == |combined|;
  assert forall j | 0 <= j < |combined| ensures result[..][j] == combined[j] {
    assert result[..][j] == result[j];
  }
  assert result[..] == combined;
  assert combined == (evens(inputs[..]) + odds(inputs[..]));
}
\end{lstlisting}
\subsection*{Program $\rightarrow$ GenAI}
\begin{lstlisting}
When using dafny resolve, the below error is emitted and resolve fails:

problem_data/p_6_2_f_move_even_to_front_array/program.dfy(82,0): Error: invalid UnaryExpression
   |
82 | }
   | ^

problem_data/p_6_2_f_move_even_to_front_array/program.dfy(81,59): Error: expected either 'by' or a semicolon following the assert expression
   |
81 |   assert combined == (evens(inputs[..]) + odds(inputs[..]));
   |                                                            ^

2 parse errors detected in program.dfy

Correct the error in code by altering only the method body. Produce and show only the Dafny body, including the curly braces that surround it. Do not show the signature nor contract. You must not use 'assume' nor 'expect' anywhere in your code.
\end{lstlisting}
\subsection*{GenAI $\rightarrow$ Program}
\textbf{System fingerprint: }null
\\\textbf{ID: }chatcmpl-Dn12d7e7YExo5a6QpFiRAyPz1vcrD
\\\textbf{Model: }gpt-5.5-2026-04-23
\\\textbf{Created at: }1780574447
\\\textbf{Finish reason: }stop
\\\textbf{Usage: }2068 tokens in, and 2602 tokens out
\begin{lstlisting}
{
  var n := inputs.Length;
  var es: seq<int> := [];
  var os: seq<int> := [];
  var i := n;

  assert inputs[n..] == [];
  assert evens(inputs[n..]) == [];
  assert odds(inputs[n..]) == [];

  while 0 < i
    invariant 0 <= i <= n
    invariant es == evens(inputs[i..])
    invariant os == odds(inputs[i..])
  {
    var oldI := i;
    var k := i - 1;
    var x := inputs[k];
    ghost var tail := inputs[oldI..];
    ghost var whole := inputs[k..];

    assert 0 <= k < n;
    assert oldI == k + 1;
    assert inputs[k..oldI] == [x];
    assert whole == inputs[k..oldI] + tail;
    assert whole == [x] + tail;
    assert |whole| != 0;
    assert whole[0] == x;
    assert whole[1..] == tail;

    if x % 2 == 0 {
      es := [x] + es;
      assert evens(whole) == [x] + evens(tail);
      assert odds(whole) == odds(tail);
      assert es == evens(inputs[k..]);
      assert os == odds(inputs[k..]);
    } else {
      os := [x] + os;
      assert evens(whole) == evens(tail);
      assert odds(whole) == [x] + odds(tail);
      assert es == evens(inputs[k..]);
      assert os == odds(inputs[k..]);
    }

    i := k;
  }

  assert i == 0;
  assert inputs[i..] == inputs[..];
  assert es == evens(inputs[..]);
  assert os == odds(inputs[..]);

  var combined := es + os;
  result := new int[|combined|](j => combined[j]);

  assert result.Length == |combined|;
  assert forall j :: 0 <= j < |combined| ==> result[j] == combined[j];
  assert |result[..]| == |combined|;
  assert forall j | 0 <= j < |combined| ensures result[..][j] == combined[j] {
    assert result[..][j] == result[j];
  }
  assert result[..] == combined;
}
\end{lstlisting}
\subsection*{Program $\rightarrow$ GenAI}
\begin{lstlisting}
When using dafny resolve, the below error is emitted and resolve fails:

problem_data/p_6_2_f_move_even_to_front_array/program.dfy(81,0): Error: invalid UnaryExpression
   |
81 | }
   | ^

problem_data/p_6_2_f_move_even_to_front_array/program.dfy(80,31): Error: expected either 'by' or a semicolon following the assert expression
   |
80 |   assert result[..] == combined;
   |                                ^

2 parse errors detected in program.dfy

Correct the error in code by altering only the method body. Produce and show only the Dafny body, including the curly braces that surround it. Do not show the signature nor contract. You must not use 'assume' nor 'expect' anywhere in your code.
\end{lstlisting}
\subsection*{GenAI $\rightarrow$ Program}
\textbf{System fingerprint: }null
\\\textbf{ID: }chatcmpl-Dn13atUTUKyIdlXphgJAH7SFL5GBo
\\\textbf{Model: }gpt-5.5-2026-04-23
\\\textbf{Created at: }1780574506
\\\textbf{Finish reason: }stop
\\\textbf{Usage: }2828 tokens in, and 3170 tokens out
\begin{lstlisting}
{
  var n := inputs.Length;
  var es: seq<int> := [];
  var os: seq<int> := [];
  var i := n;

  assert inputs[n..] == [];
  assert evens(inputs[n..]) == [];
  assert odds(inputs[n..]) == [];

  while 0 < i
    invariant 0 <= i <= n
    invariant es == evens(inputs[i..])
    invariant os == odds(inputs[i..])
  {
    var oldI := i;
    var k := i - 1;
    var x := inputs[k];
    ghost var tail := inputs[oldI..];
    ghost var whole := inputs[k..];

    assert 0 <= k < n;
    assert oldI == k + 1;
    assert inputs[k..oldI] == [x];
    assert whole == inputs[k..oldI] + tail;
    assert whole == [x] + tail;
    assert |whole| != 0;
    assert whole[0] == x;
    assert whole[1..] == tail;

    if x % 2 == 0 {
      es := [x] + es;
      assert evens(whole) == [x] + evens(tail);
      assert odds(whole) == odds(tail);
      assert es == evens(inputs[k..]);
      assert os == odds(inputs[k..]);
    } else {
      os := [x] + os;
      assert evens(whole) == evens(tail);
      assert odds(whole) == [x] + odds(tail);
      assert es == evens(inputs[k..]);
      assert os == odds(inputs[k..]);
    }

    i := k;
  }

  assert i == 0;
  assert inputs[i..] == inputs[..];
  assert es == evens(inputs[..]);
  assert os == odds(inputs[..]);

  var combined := es + os;
  result := new int[|combined|](j => combined[j]);

  assert result[..] == combined by {
    assert result.Length == |combined|;
    assert |result[..]| == |combined|;
    assert forall j :: 0 <= j < |combined| ==> result[j] == combined[j];
    assert forall j :: 0 <= j < |combined| ==> result[..][j] == combined[j];
  }
}
\end{lstlisting}
\subsection*{Program $\rightarrow$ GenAI}
\begin{lstlisting}
When using dafny verify, the below error is emitted and verify fails:

problem_data/p_6_2_f_move_even_to_front_array/program.dfy(72,45): Error: index out of range
   |
72 |   result := new int[|combined|](j => combined[j]);
   |                                              ^


Dafny program verifier finished with 3 verified, 1 error

Correct the error in code by altering only the method body. Produce and show only the Dafny body, including the curly braces that surround it. Do not show the signature nor contract. You must not use 'assume' nor 'expect' anywhere in your code.
\end{lstlisting}
\subsection*{GenAI $\rightarrow$ Program}
\textbf{System fingerprint: }null
\\\textbf{ID: }chatcmpl-Dn15FSv19Dzxlhht0mhtAIqv1z7US
\\\textbf{Model: }gpt-5.5-2026-04-23
\\\textbf{Created at: }1780574609
\\\textbf{Finish reason: }stop
\\\textbf{Usage: }3519 tokens in, and 4528 tokens out
\begin{lstlisting}
{
  var n := inputs.Length;
  var es: seq<int> := [];
  var os: seq<int> := [];
  var i := n;

  assert inputs[n..] == [];
  assert evens(inputs[n..]) == [];
  assert odds(inputs[n..]) == [];

  while 0 < i
    invariant 0 <= i <= n
    invariant es == evens(inputs[i..])
    invariant os == odds(inputs[i..])
  {
    var oldI := i;
    var k := i - 1;
    var x := inputs[k];
    ghost var tail := inputs[oldI..];
    ghost var whole := inputs[k..];

    assert 0 <= k < n;
    assert oldI == k + 1;
    assert inputs[k..oldI] == [x];
    assert whole == inputs[k..oldI] + tail;
    assert whole == [x] + tail;
    assert |whole| != 0;
    assert whole[0] == x;
    assert whole[1..] == tail;

    if x % 2 == 0 {
      es := [x] + es;
      assert evens(whole) == [x] + evens(tail);
      assert odds(whole) == odds(tail);
      assert es == evens(inputs[k..]);
      assert os == odds(inputs[k..]);
    } else {
      os := [x] + os;
      assert evens(whole) == evens(tail);
      assert odds(whole) == [x] + odds(tail);
      assert es == evens(inputs[k..]);
      assert os == odds(inputs[k..]);
    }

    i := k;
  }

  assert i == 0;
  assert inputs[i..] == inputs[..];
  assert es == evens(inputs[..]);
  assert os == odds(inputs[..]);

  var combined := es + os;
  var m := |combined|;
  result := new int[m];

  var j := 0;
  while j < m
    invariant 0 <= j <= m
    invariant result.Length == m
    invariant forall t :: 0 <= t < j ==> result[t] == combined[t]
  {
    result[j] := combined[j];
    j := j + 1;
  }

  assert j == m;
  assert forall t :: 0 <= t < m ==> result[t] == combined[t];
  assert |result[..]| == m;
  forall t | 0 <= t < m
    ensures result[..][t] == combined[t]
  {
    assert result[..][t] == result[t];
  }
  assert forall t :: 0 <= t < |result[..]| ==> result[..][t] == combined[t];
  assert result[..] == combined;
  assert combined == es + os;
  assert combined == evens(inputs[..]) + odds(inputs[..]);
}
\end{lstlisting}
\section*{Final Program}
\begin{lstlisting}
function evens(s: seq<int>): seq<int>
{
  if |s| == 0 then []
  else if s[0] % 2 == 0
       then [s[0]] + evens(s[1..])
       else evens(s[1..])
}

function odds(s: seq<int>): seq<int>
{
  if |s| == 0 then []
  else if s[0] % 2 != 0
       then [s[0]] + odds(s[1..])
       else odds(s[1..])
}

method p_6_2_f_move_even_to_front_array(inputs: array<int>) returns (result: array<int>)
	ensures result[..] == evens(inputs[..]) + odds(inputs[..])
{
  var n := inputs.Length;
  var es: seq<int> := [];
  var os: seq<int> := [];
  var i := n;

  assert inputs[n..] == [];
  assert evens(inputs[n..]) == [];
  assert odds(inputs[n..]) == [];

  while 0 < i
    invariant 0 <= i <= n
    invariant es == evens(inputs[i..])
    invariant os == odds(inputs[i..])
  {
    var oldI := i;
    var k := i - 1;
    var x := inputs[k];
    ghost var tail := inputs[oldI..];
    ghost var whole := inputs[k..];

    assert 0 <= k < n;
    assert oldI == k + 1;
    assert inputs[k..oldI] == [x];
    assert whole == inputs[k..oldI] + tail;
    assert whole == [x] + tail;
    assert |whole| != 0;
    assert whole[0] == x;
    assert whole[1..] == tail;

    if x % 2 == 0 {
      es := [x] + es;
      assert evens(whole) == [x] + evens(tail);
      assert odds(whole) == odds(tail);
      assert es == evens(inputs[k..]);
      assert os == odds(inputs[k..]);
    } else {
      os := [x] + os;
      assert evens(whole) == evens(tail);
      assert odds(whole) == [x] + odds(tail);
      assert es == evens(inputs[k..]);
      assert os == odds(inputs[k..]);
    }

    i := k;
  }

  assert i == 0;
  assert inputs[i..] == inputs[..];
  assert es == evens(inputs[..]);
  assert os == odds(inputs[..]);

  var combined := es + os;
  var m := |combined|;
  result := new int[m];

  var j := 0;
  while j < m
    invariant 0 <= j <= m
    invariant result.Length == m
    invariant forall t :: 0 <= t < j ==> result[t] == combined[t]
  {
    result[j] := combined[j];
    j := j + 1;
  }

  assert j == m;
  assert forall t :: 0 <= t < m ==> result[t] == combined[t];
  assert |result[..]| == m;
  forall t | 0 <= t < m
    ensures result[..][t] == combined[t]
  {
    assert result[..][t] == result[t];
  }
  assert forall t :: 0 <= t < |result[..]| ==> result[..][t] == combined[t];
  assert result[..] == combined;
  assert combined == es + os;
  assert combined == evens(inputs[..]) + odds(inputs[..]);
}
\end{lstlisting}
\section*{Total Token Usage}
\textbf{Input tokens: }10278
\\\textbf{Output tokens: }20066
\\\textbf{Reasoning tokens: }16755
\\\textbf{Sum of `total tokens': }30344
\section*{Experiment Timings}
\textbf{Overall Experiment} started at 1780574227701, ended at 1780574724997, lasting 497296ms (497.30 seconds)
\\\textbf{Iteration \#4} started at 1780574504515, ended at 1780574608142, lasting 103627ms (103.63 seconds)
\\\textbf{Iteration \#5} started at 1780574608160, ended at 1780574724978, lasting 116818ms (116.82 seconds)
\\\textbf{Iteration \#1} started at 1780574227702, ended at 1780574416661, lasting 188959ms (188.96 seconds)
\\\textbf{Iteration \#2} started at 1780574416662, ended at 1780574446767, lasting 30105ms (30.11 seconds)
\\\textbf{Iteration \#3} started at 1780574446768, ended at 1780574504514, lasting 57746ms (57.75 seconds)
\end{document}
